\documentclass[conference]{IEEEtran}
\usepackage{amsmath,amssymb,amsfonts}
\usepackage{graphicx}
\usepackage{booktabs}
\usepackage{cite}
\usepackage[hidelinks]{hyperref}

\begin{document}

\title{A Truncated-Laplace Incidence Model and its Moment-Matched Gamma
Approximation for the Aggregate Gain of OIRS-Assisted Optical Wireless Links}

\author{\IEEEauthorblockN{Author Name(s)}
\IEEEauthorblockA{Affiliation, Department\\
City, Country\\
Email: name@domain}}

\maketitle

\begin{abstract}
Optical intelligent reflecting surfaces (OIRSs) create controllable
non-line-of-sight (NLoS) paths in optical wireless communication (OWC)
systems, but the resulting aggregate channel gain is analytically awkward
because it depends on the \emph{cosine of the incidence angle} at each
reflecting element, a bounded random quantity. In this paper we model the
incidence cosine as a \emph{truncated-Laplace} random variable defined on a
geometry-imposed interval, derive its exact first and second moments in
closed form, and use moment matching to approximate the aggregate gain
$Z_k$ of an $M$-LED, $N$-element OIRS link by a Gamma distribution. Closed
forms for the Gamma shape $k_g$ and scale $\theta_g$ are given directly in
terms of the truncation limits. The approximation is verified against
$10^{5}$-sample Monte~Carlo simulations using the Bhattacharyya similarity
coefficient (BSC) for the probability density function (PDF) and the
Kolmogorov--Smirnov (K--S) maximum distance for the cumulative distribution
function (CDF). For the single-path reference case ($M{=}N{=}1$) the fit
already attains a BSC of 97.33\%; as the number of aggregated
paths $MN$ grows, the fit tightens monotonically toward a BSC exceeding
99.9\% and a K--S distance below $0.013$, consistent with the
central-limit behaviour of the sum. The result provides a tractable channel
model for outage and error-rate analysis of OIRS-assisted OWC.
\end{abstract}

\begin{IEEEkeywords}
Optical wireless communication, optical intelligent reflecting surface,
truncated Laplace distribution, moment matching, Gamma distribution,
channel statistics.
\end{IEEEkeywords}

\section{Introduction}
Reconfigurable/intelligent reflecting surfaces have emerged as a means of
shaping the propagation environment in wireless systems~\cite{basar2019}, and
the concept transfers naturally to the optical band, where an optical
intelligent reflecting surface (OIRS) can redirect light to establish a
controllable non-line-of-sight (NLoS) path when the direct
line-of-sight link is blocked. Optical wireless links driven by LED
sources are well studied in the line-of-sight case~\cite{komine2004}, but the
OIRS-assisted NLoS gain is harder to characterise: each reflecting element
contributes a term that scales with the \emph{cosine of the incidence angle}
$\xi$, and $\xi$ is a bounded random variable whose range is fixed by the
device geometry and the receiver field of view (FoV).

Exact distributions of such angle-dependent gains are typically intractable,
so accurate yet simple surrogate distributions are valuable for downstream
performance analysis (outage probability, bit-error rate, ergodic rate). In
this paper we (i) adopt a \emph{truncated-Laplace} model for the incidence
cosine $\xi$, which captures a sharp mode with light, bounded tails; (ii)
derive its exact first and second moments in closed form over the truncation
interval; (iii) obtain, by moment matching, a Gamma approximation of the
aggregate gain $Z_k$ of an $M$-LED, $N$-element link, with closed-form shape
and scale; and (iv) verify the approximation against Monte~Carlo simulation
using two complementary goodness-of-fit measures. The moment-matching idea
is classical~\cite{papoulis2002}; our contribution is the closed-form
truncated-Laplace moments and the demonstration that the resulting Gamma
model is asymptotically tight in the number of aggregated paths.

\emph{Notation.} $\mathbb{E}[\cdot]$ is expectation, $\mathrm{Var}[\cdot]$
variance, and $\Gamma(k_g,\theta_g)$ the Gamma distribution with shape
$k_g$ and scale $\theta_g$.

\section{Related Work}
\label{sec:related}
The idea of engineering the propagation environment with reconfigurable
surfaces was popularised for radio-frequency systems~\cite{basar2019} and
has since been carried into the optical domain, where reflectors steer
LED-generated light along controllable NLoS paths. Classical LED-based OWC
analysis, however, is dominated by the line-of-sight Lambertian
model~\cite{komine2004}, in which the received power depends deterministically
on the irradiance and incidence angles; randomness enters only through user
mobility or blockage. OIRS-assisted links differ in that the reflected
contribution of each element is modulated by the \emph{incidence-angle cosine}
at the surface, which is naturally a bounded random quantity once the user
position, panel orientation, and receiver FoV are treated statistically.

For tractability, prior OWC studies frequently replace an intractable exact
gain distribution with a moment-matched surrogate; the Gamma and log-normal
families are common choices because their moments map cleanly to physical
parameters~\cite{papoulis2002}. Our approach follows this tradition but
differs in two respects. First, we motivate the surrogate from an explicit
\emph{truncated-Laplace} model of the incidence cosine---a heavier-peaked,
strictly bounded alternative to the truncated-Gaussian---and derive its
moments in closed form. Second, rather than asserting the surrogate's
accuracy, we characterise it quantitatively with two goodness-of-fit measures
and show how the fit scales with the array size, giving the designer an
explicit validity regime.

\section{System and Channel Model}
\label{sec:model}
We consider an OWC link with $M$ LED transmitters and an OIRS of $N$
reflecting elements establishing an NLoS path to a photodetector receiver.
The aggregate NLoS channel gain is written as
\begin{equation}
Z_k \;=\; \sum_{m=1}^{M}\sum_{n=1}^{N} a_k \,\xi_{mn},
\label{eq:agg}
\end{equation}
where $\xi_{mn}$ is the cosine of the incidence angle at element $n$ for LED
$m$, and $a_k$ is a deterministic per-path scalar collecting the optical
path loss, detector responsivity/area, and element reflectivity for
configuration $k$.

\noindent\textbf{[Author to insert:]} \emph{the explicit expression for
$a_k$ from the link geometry (transmit power, Lambertian order, source--OIRS
and OIRS--receiver distances, effective areas, FoV). In the verification
below we use the worked-example value $a_k=10^{-6}$ from
Section~\ref{sec:results}, so that the statistical claims are independent of
the specific path-loss algebra; substituting the exact $a_k$ rescales
$Z_k$ but leaves the shape parameter $k_g$ and all goodness-of-fit metrics
unchanged.}

Because the $\xi_{mn}$ are bounded by the FoV and the surface geometry, and
concentrate around a nominal incidence, we model each as an independent and
identically distributed truncated-Laplace variable, defined next.

\section{Statistical Characterization of the Incidence Cosine}
\label{sec:stats}
Let $\xi$ be Laplace-distributed with location $\mu_\xi$ and scale $b_\xi$,
truncated to the geometry-imposed interval $[\xi_{\min},\xi_{\max}]$, with
$0\le\xi_{\min}$ (FoV limit) and $\xi_{\max}\le 1$ (upper geometric bound).
Introducing the centred variable $u=\xi-\mu_\xi$ on $[\alpha,\beta]$ with
\begin{equation}
\alpha=\xi_{\min}-\mu_\xi\le 0,\qquad \beta=\xi_{\max}-\mu_\xi\ge 0,
\end{equation}
the truncated density is $p(u)=\Lambda^{-1}\exp(-|u|/b_\xi)$ with
normalisation
\begin{equation}
\Lambda=\int_{\alpha}^{\beta}e^{-|u|/b_\xi}\,du
      =b_\xi\!\left(2-e^{\alpha/b_\xi}-e^{-\beta/b_\xi}\right).
\label{eq:Lambda}
\end{equation}

\subsection{First moment}
Because $|u|=-u$ on $[\alpha,0]$ and $|u|=u$ on $[0,\beta]$, the first moment
splits as
\begin{equation}
\mathbb{E}[u]=\frac{1}{\Lambda}\!\left[
\int_{\alpha}^{0}\! u\,e^{u/b_\xi}\,du
+\int_{0}^{\beta}\! u\,e^{-u/b_\xi}\,du\right].
\end{equation}
Using $\int u\,e^{u/b}\,du=e^{u/b}(bu-b^2)$ and
$\int u\,e^{-u/b}\,du=-e^{-u/b}(bu+b^2)$ and evaluating the limits, the two
constant terms at $u=0$ cancel and one obtains the closed form
\begin{equation}
\mathbb{E}[u]=\frac{b_\xi}{\Lambda}
\!\left[e^{\alpha/b_\xi}(b_\xi-\alpha)-e^{-\beta/b_\xi}(\beta+b_\xi)\right].
\label{eq:Eu}
\end{equation}

\subsection{Second moment}
Applying the same split with
$\int u^{2}e^{u/b}\,du=e^{u/b}(bu^{2}-2b^{2}u+2b^{3})$ and
$\int u^{2}e^{-u/b}\,du=-e^{-u/b}(bu^{2}+2b^{2}u+2b^{3})$, the $u=0$ terms
combine to $4b_\xi^{3}$ and
\begin{align}
\mathbb{E}[u^2]=\frac{1}{\Lambda}\Big[&\,4b_\xi^{3}
- e^{\alpha/b_\xi}\!\big(b_\xi\alpha^{2}-2b_\xi^{2}\alpha+2b_\xi^{3}\big)
\nonumber\\
&- e^{-\beta/b_\xi}\!\big(b_\xi\beta^{2}+2b_\xi^{2}\beta+2b_\xi^{3}\big)\Big].
\label{eq:Eu2}
\end{align}
The moments of the incidence cosine follow immediately:
\begin{equation}
\mathbb{E}[\xi]=\mu_\xi+\mathbb{E}[u],\qquad
\mathrm{Var}[\xi]=\mathbb{E}[u^2]-\mathbb{E}[u]^2.
\label{eq:xi}
\end{equation}
Equations \eqref{eq:Lambda}--\eqref{eq:xi} are exact for any admissible
$(\mu_\xi,b_\xi,\xi_{\min},\xi_{\max})$; no series truncation or numerical
integration is required.

\section{Moment-Matched Gamma Approximation}
\label{sec:gamma}
Since the $\xi_{mn}$ in \eqref{eq:agg} are i.i.d.\ and $a_k$ is deterministic,
the mean and variance of the aggregate gain are
\begin{equation}
\mathbb{E}[Z_k]=MN\,a_k\,\mathbb{E}[\xi],\qquad
\mathrm{Var}[Z_k]=MN\,a_k^{2}\,\mathrm{Var}[\xi].
\label{eq:Zmoments}
\end{equation}
The gain $Z_k$ is non-negative and, being a sum of many bounded positive
contributions, is unimodal and right-supported---properties shared by the
Gamma family. Matching the first two moments of a $\Gamma(k_g,\theta_g)$
variate to \eqref{eq:Zmoments} gives the closed forms
\begin{equation}
k_g=\frac{\mathbb{E}[Z_k]^2}{\mathrm{Var}[Z_k]}
   =\frac{MN\,\mathbb{E}[\xi]^2}{\mathrm{Var}[\xi]},\qquad
\theta_g=\frac{\mathrm{Var}[Z_k]}{\mathbb{E}[Z_k]}
        =\frac{a_k\,\mathrm{Var}[\xi]}{\mathbb{E}[\xi]}.
\label{eq:kgtheta}
\end{equation}
The approximate PDF and CDF of the aggregate gain are then
\begin{equation}
f_{Z_k}(z)\!\approx\!\frac{z^{k_g-1}e^{-z/\theta_g}}
{\Gamma(k_g)\,\theta_g^{k_g}},\;\;
F_{Z_k}(z)\!\approx\!\frac{\gamma\!\left(k_g,\,z/\theta_g\right)}{\Gamma(k_g)},
\label{eq:gammapdf}
\end{equation}
where $\gamma(\cdot,\cdot)$ is the lower incomplete Gamma function.

A key structural observation follows from \eqref{eq:kgtheta}: the shape
parameter scales \emph{linearly} with the number of aggregated paths,
$k_g\propto MN$, while $\theta_g$ is independent of $MN$. As $MN$ grows the
Gamma distribution therefore becomes increasingly concentrated and
Gaussian-like, matching the central-limit behaviour of the sum in
\eqref{eq:agg}. We should thus expect the approximation to be \emph{loosest}
for the single-path case $M{=}N{=}1$ and to tighten monotonically with
$MN$---a prediction we confirm numerically below.

\section{Numerical Results and Verification}
\label{sec:results}
\subsection{Setup}
We use the worked-example geometry parameters
$\mu_\xi=0.8789$, $b_\xi=0.03600$, $\xi_{\min}=0$, $\xi_{\max}=0.9574$, and
$a_k=10^{-6}$. Monte~Carlo baselines use $N_{\mathrm{mc}}=10^{5}$
realisations of $Z_k$; truncated-Laplace samples are drawn by inverse-CDF
Laplace generation followed by accept--reject rejection of samples outside
$[\xi_{\min},\xi_{\max}]$. Substituting the parameters into
\eqref{eq:xi} gives the exact single-element moments
\begin{equation}
\mathbb{E}[\xi]=0.87204,\qquad \mathrm{Var}[\xi]=1.838\times10^{-3},
\end{equation}
in close agreement with the Monte~Carlo estimates
($\widehat{\mathbb{E}}[\xi]=0.87217$,
$\widehat{\mathrm{Var}}[\xi]=1.832\times10^{-3}$), which validates the
closed-form moment expressions \eqref{eq:Eu}--\eqref{eq:Eu2}.

\subsection{Sampling methodology and complexity}
Truncated-Laplace realisations are generated by drawing an untruncated
Laplace variate---via the inverse-CDF rule
$\xi=\mu_\xi-b_\xi\,\mathrm{sgn}(U)\ln(1-2|U|)$ with $U\!\sim\!
\mathcal{U}(-\tfrac12,\tfrac12)$---and rejecting draws outside
$[\xi_{\min},\xi_{\max}]$. For the example parameters the acceptance
probability is
\begin{equation}
\Pr\{\xi_{\min}\!\le\!\xi\!\le\!\xi_{\max}\}
=F_L(\xi_{\max})-F_L(\xi_{\min})=0.9435,
\end{equation}
where $F_L$ is the untruncated Laplace CDF; i.e.\ 94.35\% of
draws are retained, so the generator is efficient and no importance
reweighting is needed. Because acceptance is high and independent across
paths, the cost of building the aggregate in \eqref{eq:agg} scales as
$\mathcal{O}(N_{\mathrm{mc}}MN)$, which is what allows the $MN$ sweep of
Table~\ref{tab:sweep} to be produced at $N_{\mathrm{mc}}=10^{5}$ in
seconds.

\subsection{Goodness-of-fit measures}
We quantify the PDF match by the Bhattacharyya similarity coefficient
\begin{equation}
\mathrm{BSC}=\sum_{i}\sqrt{P_{\mathrm{emp}}(i)\,P_{\mathrm{ana}}(i)}
\in[0,1],
\end{equation}
where $P_{\mathrm{emp}}$ and $P_{\mathrm{ana}}$ are the empirical and
analytical bin probabilities ($\mathrm{BSC}=1$ is a perfect match), and the
CDF match by the Kolmogorov--Smirnov maximum distance
\begin{equation}
D_{\mathrm{KS}}=\sup_z\bigl|F_{\mathrm{emp}}(z)-F_{Z_k}(z)\bigr|.
\end{equation}

\subsection{Single-path reference case ($M{=}N{=}1$)}
The moment-matched parameters are $k_g=413.82$ and
$\theta_g=2.107\times10^{-9}$. Fig.~\ref{fig:mn1} overlays the analytical
Gamma PDF/CDF on the Monte~Carlo histogram/empirical CDF. As anticipated in
Section~\ref{sec:gamma}, this hardest case is where the Gamma and the (left-
skewed, upper-truncated) empirical density differ most, yet the fit is
already good: $\mathrm{BSC}=97.33\%$ and
$D_{\mathrm{KS}}=0.0835$.

\begin{figure}[t]
\centering
\includegraphics[width=\columnwidth]{fit_base_MN1.png}
\caption{Verification for the single-path case $M{=}N{=}1$
($k_g=413.82$, $\theta_g=2.107\times10^{-9}$). Top: PDF; bottom: CDF.
This is the loosest case of the family
($\mathrm{BSC}=97.33\%$, $D_{\mathrm{KS}}=0.0835$).}
\label{fig:mn1}
\end{figure}

\subsection{Effect of path aggregation}
Table~\ref{tab:sweep} sweeps $MN$. Consistent with $k_g\propto MN$, the
Gamma fit tightens monotonically: the BSC rises past
99.9\% and the K--S distance falls below $0.013$.
Fig.~\ref{fig:mn16} shows the representative case $MN=16$, where the
analytical Gamma is visually indistinguishable from the simulation in both
the PDF and the CDF.

\begin{table}[t]
\centering
\caption{Gamma fit quality versus number of aggregated paths $MN$
($N_{\mathrm{mc}}=10^{5}$).}
\label{tab:sweep}
\begin{tabular}{@{}rrrr@{}}
\toprule
$MN$ & $k_g$ & BSC (\%) & $D_{\mathrm{KS}}$ \\
\midrule
1  & 413.8    & 97.37 & 0.0827 \\
2  & 827.6    & 98.94 & 0.0546 \\
4  & 1655.3   & 99.41 & 0.0379 \\
8  & 3310.6   & 99.67 & 0.0254 \\
16 & 6621.1   & 99.83 & 0.0186 \\
32 & 13242.2  & 99.90 & 0.0150 \\
64 & 26484.4  & 99.94 & 0.0127 \\
\bottomrule
\end{tabular}
\end{table}

\begin{figure}[t]
\centering
\includegraphics[width=\columnwidth]{fit_good_MN16.png}
\caption{Verification for $MN=16$
($\mathrm{BSC}=99.83\%$, $D_{\mathrm{KS}}=0.0186$). The Gamma
approximation is essentially exact once several paths are aggregated.}
\label{fig:mn16}
\end{figure}

\subsection{Application: outage probability}
The value of a closed-form CDF is that link-level metrics follow immediately.
Consider an intensity-modulation/direct-detection receiver for which a target
quality of service requires the aggregate gain to exceed a threshold
$z_{\mathrm{th}}$ (equivalently, the received electrical SNR to exceed a
target $\gamma_{\mathrm{th}}$, with $z_{\mathrm{th}}$ the corresponding gain).
Using \eqref{eq:gammapdf}, the outage probability is obtained in closed form
as
\begin{equation}
P_{\mathrm{out}}(z_{\mathrm{th}})
=\Pr\{Z_k\le z_{\mathrm{th}}\}
=F_{Z_k}(z_{\mathrm{th}})
\approx\frac{\gamma\!\left(k_g,\,z_{\mathrm{th}}/\theta_g\right)}{\Gamma(k_g)} .
\label{eq:outage}
\end{equation}
Because $k_g\propto MN$ while the per-path scale $\theta_g$ is fixed
(Section~\ref{sec:gamma}), \eqref{eq:outage} makes the array-size dependence
of the outage explicit: enlarging the OIRS increases $\mathbb{E}[Z_k]$
linearly and concentrates the distribution, sharply reducing
$P_{\mathrm{out}}$ for a fixed threshold. The same $F_{Z_k}$ can be fed to
standard averaging integrals for the average bit-error rate, so the
truncated-Laplace/Gamma pair reduces performance analysis to evaluations of
the (widely tabulated) incomplete Gamma function rather than nested numerical
integration over the incidence-angle statistics.

\subsection{Model applicability}
The truncated-Laplace assumption is appropriate when the incidence cosine is
sharply concentrated about a nominal value---as it is for a fixed or
slowly-moving receiver whose orientation varies within a limited FoV---and
strictly bounded by geometry. When the receiver orientation spreads more
uniformly over a wide FoV, a bounded-support model with a flatter body (e.g.\
truncated-Gaussian or Beta) may fit the incidence statistics better; the
moment-matching machinery of Section~\ref{sec:gamma} is agnostic to this
choice, requiring only the first two moments of $\xi$. The closed forms
\eqref{eq:Eu}--\eqref{eq:Eu2} are what make the truncated-Laplace instance
especially convenient. Validating the truncated-Laplace assumption against
ray-traced or measured incidence-angle histograms for a specific OIRS
deployment is left as future work.

\subsection{Discussion}
Two points are worth emphasising. First, the closed-form moments
\eqref{eq:Eu}--\eqref{eq:Eu2} are exact and match simulation to three
significant figures, so any residual mismatch in Figs.~\ref{fig:mn1}
and~\ref{fig:mn16} is attributable solely to the Gamma \emph{shape}
assumption, not to the moment computation. Second, because $k_g$ scales with
$MN$ while $\theta_g$ does not, a designer can read off from
Table~\ref{tab:sweep} the smallest array size for which the tractable Gamma
model is trustworthy for a target accuracy---e.g.\ $D_{\mathrm{KS}}<0.02$ is
reached by $MN\ge 16$. For very small arrays the model remains a useful but
approximate summary; for the array sizes typical of practical OIRS panels it
is effectively exact.

\section{Conclusion}
We modelled the incidence cosine of an OIRS-assisted optical wireless link as
a truncated-Laplace variable, derived its exact first and second moments in
closed form, and obtained a moment-matched Gamma approximation of the
aggregate channel gain with closed-form shape and scale. Monte~Carlo
verification via the Bhattacharyya coefficient and the Kolmogorov--Smirnov
distance confirms that the Gamma model is asymptotically tight in the number
of aggregated paths, exceeding 99.9\% similarity for moderate
array sizes. The closed-form CDF \eqref{eq:gammapdf} is directly usable in
outage-probability and error-rate analysis, which we leave to future work
together with the incorporation of the full path-loss expression for $a_k$
and measurement-based validation of the truncated-Laplace assumption.

\begin{thebibliography}{1}
\bibitem{basar2019}
E.~Basar, M.~Di~Renzo, J.~de~Rosny, M.~Debbah, M.-S.~Alouini, and R.~Zhang,
``Wireless communications through reconfigurable intelligent surfaces,''
\emph{IEEE Access}, vol.~7, pp.~116753--116773, 2019.

\bibitem{komine2004}
T.~Komine and M.~Nakagawa,
``Fundamental analysis for visible-light communication system using LED
lights,'' \emph{IEEE Trans. Consum. Electron.}, vol.~50, no.~1,
pp.~100--107, 2004.

\bibitem{papoulis2002}
A.~Papoulis and S.~U.~Pillai, \emph{Probability, Random Variables, and
Stochastic Processes}, 4th~ed. New York, NY, USA: McGraw-Hill, 2002.

\bibitem{bhattacharyya1943}
A.~Bhattacharyya, ``On a measure of divergence between two statistical
populations defined by their probability distributions,''
\emph{Bull. Calcutta Math. Soc.}, vol.~35, pp.~99--109, 1943.

\end{thebibliography}

\end{document}
